\documentclass[11pt,a4paper]{article}

\usepackage{fontspec}
\IfFontExistsTF{Times New Roman}{\setmainfont{Times New Roman}}{\setmainfont{Tinos}}
\IfFontExistsTF{Consolas}{\setmonofont{Consolas}[Scale=0.92]}{\setmonofont{Liberation Mono}[Scale=0.92]}
\usepackage[english]{babel}
\usepackage{geometry}
\usepackage{xcolor}
\usepackage{enumitem}
\usepackage{listings}
\usepackage{booktabs}
\usepackage{array}
\usepackage{amsmath,amssymb}
\usepackage{titlesec}
\usepackage{fancyhdr}
\usepackage{parskip}
\usepackage{tcolorbox}
\tcbuselibrary{skins}
\usepackage{tabularx}
\usepackage{lastpage}
\usepackage{hyperref}

\geometry{top=1.5cm,bottom=1.5cm,left=1.7cm,right=1.7cm}

\definecolor{ForestGreen}{RGB}{22,100,70}
\definecolor{SoftGreen}{RGB}{229,244,235}
\definecolor{CodeBg}{RGB}{248,250,249}
\definecolor{DarkGray}{RGB}{55,55,55}
\definecolor{primaryblue}{RGB}{22,100,70}
\definecolor{secondaryteal}{RGB}{22,100,70}
\definecolor{boxbg}{RGB}{248,250,249}
\definecolor{borderblue}{RGB}{159,198,178}

\setlength{\parindent}{0pt}
\setlength{\parskip}{0pt}
\setlist[itemize]{leftmargin=1.3em,itemsep=1pt,topsep=1pt,parsep=0pt,partopsep=0pt}
\setlist[enumerate]{leftmargin=1.5em,itemsep=1.5pt,topsep=1pt,parsep=0pt,partopsep=0pt}
\setlist[enumerate,2]{topsep=0.5pt,itemsep=0pt,parsep=0pt,partopsep=0pt}

\lstset{
  language=Python,
  basicstyle=\ttfamily\scriptsize,
  keywordstyle=\color{blue!70!black}\bfseries,
  stringstyle=\color{red!65!black},
  commentstyle=\color{ForestGreen},
  showstringspaces=false,
  columns=fullflexible,
  keepspaces=true,
  breaklines=true,
  frame=single,
  rulecolor=\color{ForestGreen!45},
  backgroundcolor=\color{CodeBg},
  xleftmargin=0.2em,
  xrightmargin=0.2em,
  aboveskip=0.2em,
  belowskip=0.2em,
  tabsize=4
}

\pagestyle{fancy}
\fancyhf{}
\setlength{\headheight}{14pt}
\fancyhead[L]{\footnotesize\textbf{DSAI1003 -- Fundamental Programming Concepts in Python}}
\fancyhead[R]{\footnotesize\textbf{Basic I/O Progress Exam -- 60 Minutes}}
\fancyfoot[L]{\scriptsize National Economics University -- College of Technology}
\fancyfoot[R]{\footnotesize Page \thepage\ of 6}
\renewcommand{\headrulewidth}{0.4pt}
\renewcommand{\footrulewidth}{0.3pt}

\titleformat{\section}{\large\bfseries\color{ForestGreen}}{}{0pt}{}
\titleformat{\subsection}{\normalsize\bfseries\color{ForestGreen}}{}{0pt}{}
\titlespacing*{\section}{0pt}{4pt}{2pt}
\titlespacing*{\subsection}{0pt}{3pt}{1.5pt}

\newtcolorbox{headerbox}{
  colback=white,
  colframe=ForestGreen,
  boxrule=0.7pt,
  arc=1.5mm,
  left=2.5mm,
  right=2.5mm,
  top=1.5mm,
  bottom=1.5mm
}

\newtcolorbox{answerbox}[2][]{
  enhanced,
  colback=white,
  colframe=DarkGray!45,
  boxrule=0.5pt,
  arc=1mm,
  left=2mm,
  right=2mm,
  top=1.2mm,
  bottom=1.2mm,
  title=\footnotesize\color{DarkGray}\textbf{#2},
  #1
}

\hypersetup{
  colorlinks=true,
  linkcolor=ForestGreen,
  urlcolor=ForestGreen,
  pdftitle={DSAI1003 - Basic Input/Output in Python Exam},
  pdfauthor={National Economics University (NEU)}
}

\begin{document}

% =========================================================================
% PAGE 1: OFFICIAL HEADER + ANSWER SHEET + SECTION A
% =========================================================================

\noindent
\begin{minipage}[t]{0.64\textwidth}
  \raggedright\small
  \textbf{NATIONAL ECONOMICS UNIVERSITY}\\
  \textbf{COLLEGE OF TECHNOLOGY}\\
  \textbf{FACULTY OF DATA SCIENCE \& ARTIFICIAL INTELLIGENCE (FDA)}
\end{minipage}%
\hfill
\begin{minipage}[t]{0.34\textwidth}
  \raggedleft\footnotesize
  \textbf{OFFICIAL MIDTERM EXAMINATION}\\
  Academic Year: 2026--2027\\
  Semester: Fall 2026 \quad Course Code: \textbf{DSAI1003}
\end{minipage}

\vspace{0.2em}

\begin{center}
  {\Large\bfseries\color{ForestGreen} FUNDAMENTAL PROGRAMMING CONCEPTS IN PYTHON}\\[0.10em]
  {\normalsize\bfseries PROGRESS EXAM: BASIC INPUT/OUTPUT IN PYTHON}\\[0.10em]
  {\normalsize\textbf{Paper Code:} \rule{1.6cm}{0.4pt} \quad | \quad \textbf{Duration: 60 Minutes} (Excluding paper distribution)}
\end{center}

\vspace{0.1em}

\begin{headerbox}
\begin{tabular*}{\textwidth}{@{\extracolsep{\fill}}ll@{}}
  \textbf{Candidate Full Name:} \dotfill & \textbf{Student ID (MSSV):} \dotfill \\[0.25em]
  \textbf{Class / Cohort:} \dotfill & \textbf{Exam Room / Seat No.:} \dotfill
\end{tabular*}

\vspace{0.25em}
\renewcommand{\arraystretch}{1.1}
\setlength{\tabcolsep}{4.5pt}
\centering
\begin{tabular}{|l|c|c|c|c||c|c|}
\hline
\textbf{Section} & \textbf{A} & \textbf{B} & \textbf{C} & \textbf{D} & \textbf{Total (100)} & \textbf{Examiner Signature} \\
\hline
\textbf{Max Score} & 20 & 25 & 25 & 30 & \textbf{100} & \\[0.35em]
\cline{1-6}
\textbf{Score} & & & & & & \\[0.45em]
\hline
\end{tabular}

\vspace{0.2em}
\raggedright
\scriptsize
\textbf{Instructions:} Closed-book examination. All electronic devices and internet access are prohibited. Answer all questions directly in the designated spaces. In programming tasks, use the \textbf{Input $\rightarrow$ Process $\rightarrow$ Output} pattern. Do \textbf{not} define user functions (\texttt{def}) or use loops/branching. Assume users enter data in the required format. Do not detach pages.
\end{headerbox}

\vspace{1mm}
\noindent
\textbf{MULTIPLE-CHOICE ANSWER SHEET} \textit{(Write A, B, C, or D in the corresponding box)}:
\begin{center}
\setlength{\tabcolsep}{5pt}
\renewcommand{\arraystretch}{1.05}
\begin{tabular}{|c|c|c|c|c|c|c|c|c|c|c|}
\hline
\textbf{Question} & \textbf{1} & \textbf{2} & \textbf{3} & \textbf{4} & \textbf{5} & \textbf{6} & \textbf{7} & \textbf{8} & \textbf{9} & \textbf{10} \\
\hline
\textbf{Answer} & & & & & & & & & & \\[2mm]
\hline
\end{tabular}
\end{center}
\vspace{-2mm}
\section*{Section A: Basic Input/Output Concepts \hfill \normalsize\normalfont(20 points)}
\textit{Answer all 10 questions. Each correct answer carries 2 points. Suggested time: 12 minutes.}

\begin{enumerate}[label=\textbf{Q\arabic*.},leftmargin=2em,itemsep=2pt,topsep=1pt]

\item Which built-in Python function is primarily used to display data on the screen?
\begin{enumerate}[label=\Alph*.,leftmargin=2em]
  \item \texttt{input()}
  \item \texttt{print()}
  \item \texttt{type()}
  \item \texttt{float()}
\end{enumerate}

\item What is printed by \texttt{print(10, 20)} using the default behavior of \texttt{print()}?
\begin{enumerate}[label=\Alph*.,leftmargin=2em]
  \item \texttt{1020}
  \item \texttt{10,20}
  \item \texttt{10 20}
  \item \texttt{10:20}
\end{enumerate}

\item What is the purpose of the \texttt{sep} parameter in \texttt{print()}?
\begin{enumerate}[label=\Alph*.,leftmargin=2em]
  \item It changes the data type of printed values.
  \item It controls the separator placed between multiple printed values.
  \item It controls whether the user may enter data.
  \item It rounds floating-point values automatically.
\end{enumerate}

\item What is the exact output of \texttt{print("A", "B", "C", sep="-")}?
\begin{enumerate}[label=\Alph*.,leftmargin=2em]
  \item \texttt{A B C}
  \item \texttt{A-B-C}
  \item \texttt{ABC-}
  \item \texttt{A,B,C}
\end{enumerate}

\item What is the primary effect of specifying \texttt{end=" "} in a \texttt{print()} statement?
\begin{enumerate}[label=\Alph*.,leftmargin=2em]
  \item It places a space after the output instead of the default newline.
  \item It places spaces between all arguments.
  \item It terminates the Python program.
  \item It converts the output into a string.
\end{enumerate}

\newpage

% ==============================================================================
% PAGE 2: SECTION A, Q6--Q10 + SECTION B Q11
% ==============================================================================

\item What data type does \texttt{input()} return, even when the user types digits such as \texttt{18}?
\begin{enumerate}[label=\Alph*.,leftmargin=2em]
  \item \texttt{int}
  \item \texttt{float}
  \item \texttt{str}
  \item The type is selected automatically from the typed value.
\end{enumerate}

\item Suppose the user enters \texttt{10} and then \texttt{20}. What does the following program print?
\begin{lstlisting}
x = input("x = ")
y = input("y = ")
print(x + y)
\end{lstlisting}
\begin{enumerate}[label=\Alph*.,leftmargin=2em]
  \item \texttt{30}
  \item \texttt{1020}
  \item \texttt{10 20}
  \item An error occurs before input is read.
\end{enumerate}

\item Which statement correctly reads a decimal price from the keyboard for use in arithmetic?
\begin{enumerate}[label=\Alph*.,leftmargin=2em]
  \item \texttt{price = input(float("Enter price: "))}
  \item \texttt{price = float(input("Enter price: "))}
  \item \texttt{price = int("Enter price: ")}
  \item \texttt{price = print(input("Enter price: "))}
\end{enumerate}

\item What does \texttt{:.2f} mean in an f-string such as \texttt{f"Price: \{price:.2f\}"}?
\begin{enumerate}[label=\Alph*.,leftmargin=2em]
  \item Convert the value to an integer with two digits.
  \item Display the floating-point value with two digits after the decimal point.
  \item Repeat the value twice.
  \item Add two to the value before displaying it.
\end{enumerate}

\item If \texttt{x = 10}, what is printed by \texttt{print("x")}?
\begin{enumerate}[label=\Alph*.,leftmargin=2em]
  \item \texttt{10}
  \item \texttt{x}
  \item \texttt{"x"}
  \item \texttt{int}
\end{enumerate}

\end{enumerate}

\section*{Section B: Output Tracing, Type Conversion \& Error Diagnosis \hfill \normalsize\normalfont(25 points)}
\textit{Answer all 5 questions. Each question carries 5 points. Suggested time: 15 minutes.}

\subsection*{Q11. Exact Output with \texttt{sep} and \texttt{end} (5 points)}
Predict the \textbf{exact screen output}, including spaces and line breaks:

\begin{lstlisting}
print("Python", "I/O", sep="::", end=" | ")
print("Input", "Process", "Output", sep=" -> ")
print("Done")
\end{lstlisting}

\begin{answerbox}[height=3.0cm]{Output Box for Q11}
\end{answerbox}

\newpage

% ==============================================================================
% PAGE 3: SECTION B, Q12--Q15
% ==============================================================================

\subsection*{Q12. Strings versus Numbers after \texttt{input()} (5 points)}
Assume the user enters \texttt{5} for each prompt. Determine the exact output of each part.

\begin{lstlisting}
a = input("a = ")
print(a * 3)

b = int(input("b = "))
print(b * 3)
\end{lstlisting}

\begin{itemize}[leftmargin=2em]
  \item Output of \texttt{print(a * 3)}: \underline{\hspace{5cm}}
  \item Output of \texttt{print(b * 3)}: \underline{\hspace{5cm}}
  \item Briefly explain why the two outputs differ: \hrulefill
\end{itemize}

\subsection*{Q13. Diagnose and Correct Common Input Errors (5 points)}
For each case, state the problem and write a corrected statement where appropriate.

\renewcommand{\arraystretch}{1.35}
\begin{tabularx}{\textwidth}{|p{4.9cm}|X|X|}
\hline
\textbf{Code / Situation} & \textbf{Problem} & \textbf{Correction / Result} \\
\hline
\texttt{age = input("Age: ")}\\
\texttt{next\_age = age + 1} & & \\[4mm]
\hline
\texttt{age = int(input("Age: "))}\\
User enters \texttt{eighteen} & & \\[4mm]
\hline
\texttt{x = 10}\\
\texttt{print("x")}\\
Goal: display the value of \texttt{x} & & \\[4mm]
\hline
\end{tabularx}

\subsection*{Q14. f-String Tracing and Formatting (5 points)}
Predict the exact output:

\begin{lstlisting}
product = "Notebook"
price = 15.5
quantity = 3
total = price * quantity
print(f"Product: {product}")
print(f"Quantity: {quantity}")
print(f"Total: {total:.2f}")
\end{lstlisting}

\begin{answerbox}[height=2.35cm]{Output Box for Q14}
\end{answerbox}

\subsection*{Q15. Identify Input, Process, and Output (5 points)}
Consider the program:

\begin{lstlisting}
celsius = float(input("Celsius: "))
fahrenheit = 9 / 5 * celsius + 32
print(f"Fahrenheit: {fahrenheit:.2f}")
\end{lstlisting}

Identify the corresponding components of the Input $\rightarrow$ Process $\rightarrow$ Output model:
\begin{itemize}[leftmargin=2em]
  \item \textbf{Input:} \hrulefill
  \item \textbf{Process:} \hrulefill
  \item \textbf{Output:} \hrulefill
\end{itemize}

\newpage

% ==============================================================================
% PAGE 4: SECTION C
% ==============================================================================

\section*{Section C: Program Design Using the Input $\rightarrow$ Process $\rightarrow$ Output Model \hfill \normalsize\normalfont(25 points)}
\textit{Answer all tasks. Suggested time: 15 minutes.}

\begin{tcolorbox}[colback=boxbg,colframe=secondaryteal,arc=1.5mm,boxrule=0.6pt]
\textbf{Problem Specification: Student Score Average Report}\\[1mm]
A program receives three floating-point scores from the user, calculates their arithmetic average, and displays the result with exactly two digits after the decimal point.
\end{tcolorbox}

\vspace{1mm}
\textbf{Task 1: Identify the IPO Components} (6 points)
\begin{itemize}[leftmargin=2em]
  \item \textbf{Input(s):} \hrulefill\\[4mm]
  \item \textbf{Process:} \hrulefill\\[4mm]
  \item \textbf{Output:} \hrulefill
\end{itemize}

\vspace{1mm}
\textbf{Task 2: Write the Program Logic in Ordered Steps} (7 points)\\
Write clear sequential steps using the Input $\rightarrow$ Process $\rightarrow$ Output structure. Your steps should state the required type conversion and output formatting.

\begin{answerbox}[height=8.0cm]{Program-Logic Work Area}
\end{answerbox}

\vspace{1mm}
\textbf{Task 3: Hand Trace} (6 points)\\
For scores \textbf{7.5}, \textbf{8.0}, and \textbf{9.0}, show the calculation and the exact formatted result.

\begin{answerbox}[height=3.7cm]{Hand-Trace Calculation}
\end{answerbox}

\vspace{1mm}
\textbf{Task 4: Explain Two Design Choices} (6 points)
\begin{enumerate}[label=\alph*),leftmargin=2em]
  \item Why should each score be converted with \texttt{float()} rather than used directly from \texttt{input()}?\\[5mm]
  \hrulefill
  \item Why is \texttt{:.2f} useful when displaying the average score?\\[5mm]
  \hrulefill
\end{enumerate}

\newpage

% ==============================================================================
% PAGE 5: SECTION D, PROBLEM 1
% ==============================================================================

\section*{Section D: Applied Programming in Python \hfill \normalsize\normalfont(30 points)}
\textit{Write complete, syntactically correct Python statements using the \textbf{Input $\rightarrow$ Process $\rightarrow$ Output} pattern. Do \textbf{not} define custom functions (\texttt{def}), use loops, or use branching. Suggested time: 18 minutes.}

\subsection*{Problem 1: Purchase Invoice Calculator (15 points)}
Write a Python program that creates a simple purchase summary.

\textbf{Requirements:}
\begin{enumerate}[leftmargin=2em,label=\arabic*.,itemsep=2pt]
  \item Read the \texttt{product\_name} using \texttt{input()}.
  \item Read the \texttt{unit\_price} and convert it to \texttt{float}.
  \item Read the \texttt{quantity} and convert it to \texttt{int}.
  \item Compute \texttt{total = unit\_price * quantity}.
  \item Display an invoice using f-strings. Show both \texttt{unit\_price} and \texttt{total} with exactly \textbf{2 digits after the decimal point}.
\end{enumerate}

The output should follow this general form:
\begin{lstlisting}
Product: Notebook
Unit price: 15.50
Quantity: 3
Total: 46.50
\end{lstlisting}

\begin{answerbox}[height=14.7cm]{Python Solution for Problem 1}
\end{answerbox}

\newpage

% ==============================================================================
% PAGE 6: SECTION D, PROBLEM 2
% ==============================================================================

\subsection*{Problem 2: Convert Seconds into Minutes and Remaining Seconds (15 points)}
Write a Python program that converts a total number of seconds into full minutes and remaining seconds.

\textbf{Requirements:}
\begin{enumerate}[leftmargin=2em,label=\arabic*.,itemsep=2pt]
  \item Prompt the user for \texttt{total\_seconds} and convert the input to \texttt{int}.
  \item Compute the number of full minutes using integer floor division (\texttt{//}).
  \item Compute the remaining seconds using the remainder operator (\texttt{\%}).
  \item Display the original number of seconds, the full minutes, and the remaining seconds clearly.
\end{enumerate}

For example, if the user enters \texttt{135}, the program should produce values equivalent to:
\begin{lstlisting}
Seconds: 135
Minutes: 2
Remaining seconds: 15
\end{lstlisting}

\begin{answerbox}[height=15.4cm]{Python Solution for Problem 2}
\end{answerbox}

\vfill
\begin{center}
  \rule{0.5\textwidth}{0.4pt}\\[1mm]
  \textbf{\color{primaryblue}--- END OF EXAM PAPER ---}
\end{center}

\end{document}
